\documentclass{situationist-preprint}

\title{Sortition and Deliberative Quality}
\msid{POL-14-b}
\msfamily{POL-14}
\msdiscipline{Political Philosophy}
\mstype{Companion manuscript}
\mskeywords{sortition; deliberation; citizens' assemblies; cognitive diversity; polarization}

\begin{document}
\maketitle

\begin{abstract}
POL-14-a argues that at least one legislative chamber should be chosen by lot, on the ground that sortition improves descriptive representativeness. It is often objected that an allotted chamber would deliberate badly: its members lack expertise, experience and the discipline of party. We argue that the objection has the comparison the wrong way round. Deliberative quality depends on the range of perspectives present, on members' freedom to change their minds, and on the information available to them. Sortition does better than election on the first two and is equal on the third when the chamber is properly supported. Evidence from deliberative polls and citizens' assemblies supports this. An allotted chamber deliberates at least as well as an elected one, and better in the respects that bear on the quality of legislation.
\end{abstract}

\section{The objection}

The case of POL-14-a is a case about who is present in the legislature. A natural objection is that presence is not enough: legislation is the product of deliberation, and the qualities that election selects for, education, articulacy and political experience, are the qualities deliberation requires. On this view an allotted chamber trades the quality of deliberation for the fidelity of a sample, and the trade is a bad one.

We take deliberation to be the weighing of reasons bearing on a decision, by a group whose members can be moved by those reasons. Its quality, for legislative purposes, has three components.

\begin{definition}\label{d:quality}
A legislative body deliberates well to the extent that (Q1) the considerations bearing on its decisions are brought before it, (Q2) its members' positions respond to the considerations so brought, and (Q3) its members have access to the information needed to assess them.
\end{definition}

\section{The comparison}

\begin{proposition}\label{p:breadth}
Other things equal, an allotted chamber does better than an elected one on Q1.
\end{proposition}

\begin{proof}
By Principle~P1 of POL-14-a, which considerations a body notices depends in part on the experiences and social positions of its members. By Lemmas~3.1 and~3.2 of POL-14-a, election produces bodies whose positions diverge systematically from the population's, and lot produces bodies whose positions match it within small sampling error. An allotted chamber therefore brings before itself a range of considerations closer to the range that bears on legislation for the whole population.
\end{proof}

The proposition has a further support in work on cognitive diversity. \citet{HongPage2004} show, in a model of problem solving, that under certain conditions a randomly chosen group of capable agents can outperform a group of the individually best performers, because the latter tend to share heuristics and to get stuck at the same points. \citet{Landemore2013} uses the result to argue that inclusive assemblies have an epistemic advantage over selected ones. The model's conditions are idealized, and we do not rest the argument on it; but it shows that individual ability is not the only, or the main, determinant of group performance.

\begin{proposition}\label{p:responsiveness}
Other things equal, an allotted chamber does better than an elected one on Q2.
\end{proposition}

\begin{proof}
An elected legislator's position on a measure is constrained by her party's program, by the commitments on which she was elected and by the prospect of re-election. These constraints are what make an elected chamber accountable, and Principle~P2 of POL-14-a requires that one chamber have them. But they also mean that her position does not respond to the considerations raised in debate except within narrow limits. A member chosen by lot has no program to keep, no party whip and no re-election to seek, and her position can follow the reasons \citep{Guerrero2014}.
\end{proof}

Q3 depends on the support a chamber receives, not on how its members are selected. Elected legislators have staff and committees; an allotted chamber needs the same, together with time, facilitation and access to expert testimony from more than one side. With such support, the evidence is that ordinary citizens deliberate competently on complex questions. Participants in deliberative polls gain substantially in factual knowledge and their opinions change in ways that track the information they receive \citep{Fishkin2009}. The Irish Citizens' Assembly of 2016--2018 considered, among other matters, the constitutional provision on abortion, heard evidence from medical, legal and ethical experts and from advocates on both sides, and produced recommendations that were taken up in a subsequent referendum \citep{FarrellSuiterHarris2019}.

\begin{theorem}\label{t:main}
An allotted chamber, adequately supported, deliberates at least as well as an elected chamber, and better with respect to Q1 and Q2.
\end{theorem}

\begin{proof}
By Propositions~\ref{p:breadth} and~\ref{p:responsiveness}, and the observation that Q3 is a matter of support available to either kind of chamber.
\end{proof}

\section{Objections}

\begin{objection}
Groups of like-minded people deliberating together become more extreme, not more reasonable.
\end{objection}

\begin{reply}
The law of group polarization concerns groups whose members begin with a shared inclination \citep{Sunstein2002}. A large random sample is the kind of group least likely to share one, because it reproduces the distribution of opinion in the population. Polarization is more to be feared in bodies selected for agreement, such as party caucuses, than in bodies selected at random.
\end{reply}

\begin{objection}
An allotted chamber dependent on staff and experts for information will be captured by them.
\end{objection}

\begin{reply}
The risk is real, and it is not peculiar to sortition: elected legislators also depend on staff and lobbyists. It is reduced by rules requiring that evidence be heard from more than one side, by rotating membership, which prevents the formation of stable relationships of dependence, and by publicity of the evidence heard. These are features of design, and they are part of what \enquote{adequately supported} in Theorem~\ref{t:main} requires.
\end{reply}

\begin{objection}
Citizens' assemblies have advised; they have not legislated. The evidence does not show how an allotted chamber with binding power would behave.
\end{objection}

\begin{reply}
This is correct, and it limits the inductive support for Theorem~\ref{t:main}. But the arguments for Propositions~\ref{p:breadth} and~\ref{p:responsiveness} do not depend on the advisory status of the bodies studied: they turn on who is present and on what constrains members' positions. Binding power would raise the stakes of capture, which is a reason for the safeguards just described, and would leave both arguments intact.
\end{reply}

\section{Conclusion}

The objection from deliberative quality does not tell against sortition. An allotted chamber brings more of the relevant considerations before itself and is freer to follow them; given the information that any legislature needs, it deliberates at least as well as an elected chamber. This reinforces the conclusion of POL-14: at least one legislative chamber should be chosen by lot, and the representativeness that lot secures also improves deliberation.

\begin{thebibliography}{9}\small
\bibitem[Farrell et~al.(2019)]{FarrellSuiterHarris2019} Farrell, D. M., Suiter, J. and Harris, C. (2019). \enquote{Systematizing} constitutional deliberation: the 2016--18 citizens' assembly in Ireland. \emph{Irish Political Studies}, 34(1).
\bibitem[Fishkin(2009)]{Fishkin2009} Fishkin, J. S. (2009). \emph{When the People Speak: Deliberative Democracy and Public Consultation}. Oxford University Press.
\bibitem[Guerrero(2014)]{Guerrero2014} Guerrero, A. A. (2014). Against elections: the lottocratic alternative. \emph{Philosophy \& Public Affairs}, 42(2), 135--178.
\bibitem[Hong and Page(2004)]{HongPage2004} Hong, L. and Page, S. E. (2004). Groups of diverse problem solvers can outperform groups of high-ability problem solvers. \emph{Proceedings of the National Academy of Sciences}, 101(46), 16385--16389.
\bibitem[Landemore(2013)]{Landemore2013} Landemore, H. (2013). \emph{Democratic Reason: Politics, Collective Intelligence, and the Rule of the Many}. Princeton University Press.
\bibitem[Sunstein(2002)]{Sunstein2002} Sunstein, C. R. (2002). The law of group polarization. \emph{Journal of Political Philosophy}, 10(2), 175--195.
\end{thebibliography}

\end{document}
